\section{Predictive Robustness}
\label{sec:predictive_robustness_results}

The selected market-only XGBoost specification is re-estimated for the alternative downside outcomes defined in Section~\ref{sec:robustness}. Table~\ref{tab:predictive-robustness} reports held-out performance for downside volatility, the worst five-day loss, and maximum drawdown over 21 and 126 trading days, together with the primary 63-day target.

\begin{table}[!htbp]
\centering
\footnotesize
\caption{Held-out predictive robustness of the selected XGBoost specification.}
\label{tab:predictive-robustness}
\begin{tabular}{lrrrrrr}
\toprule
Target & Quarters & $\mathrm{MAE}$ & $\mathrm{RMSE}$ & $R^2$ & $\overline{\rho}$ & D10--D1 spread \\
\midrule
63-day maximum drawdown & 8 & 0.0912 & 0.1258 & 0.497 & 0.738 & 41.6\% \\
63-day downside volatility & 8 & 0.1211 & 0.2362 & 0.482 & 0.840 & 76.5\% \\
Worst five-day loss & 8 & 0.0594 & 0.0942 & 0.437 & 0.748 & 27.9\% \\
21-day maximum drawdown & 8 & 0.0627 & 0.0930 & 0.394 & 0.690 & 25.2\% \\
126-day maximum drawdown & 7 & 0.1113 & 0.1455 & 0.539 & 0.752 & 52.4\% \\
\bottomrule
\end{tabular}
\end{table}


\noindent
The absolute MAE values are not directly comparable because the targets have different scales. Rank correlations remain strong across all definitions, ranging from approximately 0.69 to 0.84. Positive Decile~10--Decile~1 spreads further indicate meaningful separation, including for the worst five-day loss and 21-day maximum drawdown. The strongest results occur for 63-day downside volatility, although its different scale and interpretation preclude direct comparison with drawdown outcomes. The predicted-risk decile profiles and quarterly diagnostics in Appendix~\ref{app:predictive_robustness_stability} show that ranking performance remains positive and comparatively stable across outcomes and individual test quarters. They also compare quarterly MAE with each target's overall test MAE, identifying relatively difficult quarters without making invalid comparisons between differently scaled outcomes.

The feature-importance comparison in Figure~\ref{fig:predictive-robustness-importance} further shows that recent downside volatility, past drawdown, and market capitalization remain the leading inputs across target definitions, although their relative weights change with the horizon and outcome.

\begin{figure}[!htbp]
    \centering
    \includegraphics[width=0.88\textwidth]{04_10_xgboost_robustness_feature_importance_heatmap.pdf}
    \caption[Feature-importance stability across robustness targets]{SHAP importance shares across the primary and alternative downside-risk targets.}
    \label{fig:predictive-robustness-importance}
\end{figure}
